\documentclass[11pt,a4paper]{article}
\usepackage[margin=18mm]{geometry}
\usepackage{amsmath,booktabs,array,tabularx}
\usepackage{cite,microtype}
\renewcommand{\citepunct}{], [}
\renewcommand{\citedash}{]--[}
\input{glyphtounicode}
\pdfgentounicode=1
\usepackage[hidelinks]{hyperref}
\hypersetup{pdftitle={An established traffic backbone for an elementary VSL task},
 pdfauthor={Differentiable traffic task review},
 pdfsubject={Literature comparison and formed modeling backbone; revised draft}}
\setlength{\parindent}{0pt}
\setlength{\parskip}{4pt}
\setlength{\abovedisplayskip}{4pt}
\setlength{\belowdisplayskip}{4pt}
\setlength{\abovedisplayshortskip}{3pt}
\setlength{\belowdisplayshortskip}{3pt}
\setlength{\tabcolsep}{4pt}
\renewcommand{\arraystretch}{1.10}
\newcommand{\head}[1]{\par\vspace{5pt}\noindent\textbf{#1}\par\nobreak\vspace{1pt}}
\newenvironment{points}{\begin{list}{$\bullet$}{%
 \setlength{\leftmargin}{1.3em}\setlength{\labelwidth}{0.6em}%
 \setlength{\labelsep}{0.4em}\setlength{\itemsep}{2pt}%
 \setlength{\parsep}{0pt}\setlength{\topsep}{2pt}}}{\end{list}}
\begin{document}
\begin{center}
{\large\bfseries An established traffic backbone for an elementary VSL task}\\[3pt]
Literature comparison and formed modeling backbone --- revised draft, 21 October 2026
\end{center}
\vspace{-3pt}

\head{1. What a suitable backbone must provide}
Can an established variable speed limit (VSL) model support the information,
control and gradient questions with few traffic-model changes? Consider an
ending lane E, its adjacent continuing lane A and farther lane B. Initial
states can share the same E and combined A+B traffic at every road position
while differing in A/B loading. We ask whether that difference changes merge
service, useful common-speed advice and its sensitivity. A suitable backbone
must represent the distinction, conserve traffic, respect receiving space
and count every initial unit's time until exit. Simplicity concerns the
traffic rules we must author and justify, not citation count.

\head{2. What the closest published models offer}
Aggregate VSL models support questions about traffic location or history;
lane models additionally resolve A/B loading. The inspected starting points
offer the following tradeoff.

\begin{center}
\textbf{Table 1. Published starting points and their fit to the question}\par
\vspace{3pt}
\begin{tabularx}{\linewidth}{@{}>{\raggedright\arraybackslash}p{0.19\linewidth}
 >{\raggedright\arraybackslash}p{0.37\linewidth}
 >{\raggedright\arraybackslash}X@{}}
\toprule
Published model & What is already supplied & Direct finding for this review\\
\midrule
Jin and Jin~\cite{jinjin}
& Spatial and link-average traffic; VSL entry control, capacity drop and
upstream holding queues.
& Simpler if the question concerns \emph{where} traffic lies along the
road. It changes the A/B information question.\\[3pt]
Wei, Han and Wang~\cite{wei}
& Link traffic retaining entry/speed history, interacting speed cohorts
and capacity drop.
& A strong aggregate alternative, but no explicit A/B loading. Native
advice acts on entering traffic.\\[3pt]
Nagalur Subraveti et al.~\cite{ns}
& Successive lane-cell densities, lane-change incentives and a lane-drop
case; Han congestion lineage~\cite{han16}.
& Selected lane backbone. Adopt Han's VSL sending restriction~\cite{han17}
and explicitly replace movement realization.\\[3pt]
Carey et al.~\cite{carey}
& Multilane cell transmission, mandatory movements and receiving-allocation
alternatives.
& Supports conservative service and forward fallback. More VSL/lane-end
task closure would be needed as the main backbone.\\
\bottomrule
\end{tabularx}
\end{center}

\textbf{Selection.} Use the Subraveti lane-density framework and its Han
congestion/VSL lineage, with one explicit conservative movement replacement.
This retains the chosen lane-information question. Jin--Jin remains a
genuine simpler route if the question changes to spatial information.
The comparison concerns inspected methods, not a globally best-model claim;
the exact cohort/control scenario and movement composition below are ours.

\head{3. The formed modeling backbone}
\textbf{3.1. Physical picture, task and root assumptions.}
This is a \emph{lane-resolved cell-transmission-type (CTM) model} using
Han's extended first-order budgets. Physical cells retain density and finite
storage throughout the road; there is no separate link transmission model
(LTM) link. E can change into A at upstream cell boundaries; traffic in the
final E cell must merge at the exit.
Unserved mass stays in its current cell, so congestion can extend upstream.
The causal chain is:
\begin{center}
\fbox{Lane stocks}\ $\longrightarrow$\ \fbox{Desires and budgets}\
$\longrightarrow$\ \fbox{Shared movements}\ $\longrightarrow$\ \fbox{New stocks}
\end{center}
Use $N=10$ sections, $\Delta x=0.1$\,km and $\Delta t=1/600$\,h (6\,s).
At step $k$, density $\rho^k_{i\ell}$ and continuous stock
$n^k_{i\ell}=\Delta x\rho^k_{i\ell}$ describe section $i$, lane
$\ell\in\{E,A,B\}$. Initialize a finite cohort independently of constant
common advice $v\in[30,60]$\,km/h, with no later arrivals. Free A/B outlets
complete trips. Assume common lane flow--density parameters, fixed geometry,
old traffic/space for every movement, equal positive priorities and no extra
merging-headway multiplier. The step satisfies $u_0\Delta t,w\Delta t\leq\Delta x$;
$u_0$ is nominal free speed and $w$ the backward-wave speed magnitude.

\textbf{3.2. How much traffic can move or be received?}
The full Han budgets~\cite{han16,han17} give a sending rate $D$ and a
shared receiving rate $R$. The discharge ceiling $Q(x)$ equals nominal
capacity $C$ below critical density $\rho_c$, then falls linearly to
$(1-\delta)C$ at jam density $k_J$:
\begin{align}
 C&=\frac{u_0wk_J}{u_0+w},\quad \rho_c=C/u_0,\quad
 Q(x)=\min\!\left\{C,C\left[1-\delta\frac{x-\rho_c}{k_J-\rho_c}\right]\right\},\label{eq:ceiling}\\
 D_{i\ell}&=\min\{v\rho_{i\ell},Q(\rho_{i-1,\ell})\},\quad
 R(x,y)=\min\{Q(x),w(k_J-y),w(k_J-x)+w_d(x-y)\}.\label{eq:budgets}
\end{align}
Here $x,y$ are the current upstream/downstream densities in the receiving
lane. The three receiving bounds express discharge capacity, downstream
space and Han's additional congestion-wave constraint. Its fixed recovery
slope is $w_d=Q_J/(k_J-Q_J/u_0)$, $Q_J=(1-\delta)C$.
The member uses $u_0=60$, $w=20$\,km/h, $k_J=150$\,veh/km/lane and
$\delta=0.1$; initial densities lie in $[0,k_J]$.
Available amounts are $d_{i\ell}=\Delta t D_{i\ell}$ and
$r_{ij}=\Delta t R(\rho_{ij},\rho_{i+1,j})$: all straight and diagonal
arrivals to target $j$ share that \emph{one} $r_{ij}$.
At the outlets use $r_{Nj}=\Delta t R(\rho_{Nj},0)$, $j=A,B$;
upstream boundary densities are zero with no inflow. Advice caps sending;
receiving changes through the evolving densities, not directly through $v$.

\textbf{3.3. Why and where does early merging occur?}
The retained Subraveti incentive framework~\cite{ns} compares perceived
lane densities $K$, using current and two downstream cells with weights
$2{:}2{:}1$~\cite{thesis}. For a nonempty origin and adjacent destination,
\begin{equation}
 P_{i,\ell j}=\left[\frac{I_{i,\ell j}K_{i\ell}-K_{ij}}
 {K_{i\ell}+K_{ij}}\right]_+,\qquad
 I=1+I^r+I^{kr}+I^{coop},\qquad [x]_+=\max(x,0).\label{eq:desire}
\end{equation}
Less crowded destinations attract traffic. Within the last $D_r=0.75$\,km,
E-to-A route bias grows as $I^r=(1-g_i/D_r)^3$, where
$g_i=(N-i)\Delta x$; it is zero farther upstream. This adds urgency, not
a fixed merge fraction or a hard start line: density incentives can cause
earlier E-to-A movement. Keep-right and cooperation use the member's declared
mode/direction conventions; free A can be encouraged toward B to
make room. Missing downstream E cells are omitted and weights renormalized;
A/B lookahead extends with zero outlet density. Empty origins make no
requests. E, A and B name current lanes, so early-merged traffic becomes A.

\textbf{3.4. Which attempts succeed, and what happens at the choke?}
Our replacement of native same-cell lateral/adjusted-forward coupling
separates desire from service. Requests use existing forward or adjacent
cells/outlets. At one boundary, temporarily omit $i$.
Normalize lateral desires, leave the rest for straight
travel, and form requests $m$ from each donor's one sending amount:
\begin{equation}
 b_{\ell j}=\frac{P_{\ell j}}{\max(1,\sum_{h\ne\ell}P_{\ell h})}\ (j\ne\ell),
 \quad b_{\ell\ell}=1-\sum_{j\ne\ell}b_{\ell j},\quad
 m_{\ell j}=d_\ell b_{\ell j}.\label{eq:requests}
\end{equation}
Adapt Carey receiving/fallback principles and Wright priority sharing
\cite{carey,wright}: each target serves requests at a common capped level,
\begin{equation}
 z_{\ell j}=\min(m_{\ell j},\lambda_j),\qquad
 \sum_\ell z_{\ell j}=\min\!\left(r_j,\sum_\ell m_{\ell j}\right).\label{eq:sharing}
\end{equation}
Requests below the shared level are fully served; unused entitlement returns to the others.
Residual donor traffic gets one straight fallback $y_\ell$, the minimum of
its unused sending and its own lane's unused receiving amount. Thus
$f_{\ell j}=z_{\ell j}+\mathbf1_{\{j=\ell\}}y_\ell$ is realized movement.
An unsuccessful E-to-A attempt can therefore continue in E if space remains.
At final E override $b_{EA}=1$ and $y_E=0$: E must use A's outlet or wait.
It competes with A and any B-to-A request for the same outlet budget.
Equal priorities are not a constant realized merge proportion; demand and
receiving state determine the flow. E's use leaves less service for the
other requests, whose retained stock can restrict later receiving.
There is no separate speed-mismatch or acceleration-recovery meter.

Every movement advances one section, diagonally when changing lanes.
Apply all flows from the old state simultaneously:
\begin{equation}
 n^{k+1}_{i\ell}=n^k_{i\ell}-\sum_j f^k_{i,\ell j}
                             +\sum_h f^k_{i-1,h\ell}.\label{eq:update}
\end{equation}
Set $f_0=0$; outlet flows $f_N$ complete trips. Matched additions/subtractions
conserve mass, while donor/receiver limits protect stock and storage.

\textbf{3.5. Complete objective, information comparison and sensitivity.}
Minimize total system time $J$, counting all initial mass until exit.
Once remaining mass is below
$\Delta x\rho_c$, all requests receive service and a unit in section $i$
has remaining mean time $(N-i+1)\Delta x/v$. Choose a fixed index $K_*$
inside that regime. With $h_{i\ell}=N-i+1$ and $\mathbf1^\top n^k$ the
total unfinished mass, the exact prefix-plus-residual objective is
\begin{equation}
 J(v)=\Delta t\sum_{k<K_*}\mathbf1^\top n^k
             +\frac{\Delta x}{v}h^\top n^{K_*}.\label{eq:objective}
\end{equation}
The residual counts future time rather than deleting small remaining stock.
Write the complete update as $n^{k+1}=F(n^k,v)$. Between branch switches,
$s^0=0$, $s^{k+1}=F_n^{[k]}s^k+F_v^{[k]}$ propagates
$s^k=\partial n^k/\partial v$ through budgets, desires, allocation and stock.
Then
\begin{equation}
 J'(v)=\Delta t\sum_{k<K_*}\mathbf1^\top s^k
       +\frac{\Delta x}{v}h^\top s^{K_*}
       -\frac{\Delta x}{v^2}h^\top n^{K_*}.\label{eq:sensitivity}
\end{equation}
These terms measure earlier traffic changes, residual-stock changes and
direct residual travel time. Hold a valid $K_*$ fixed locally; at native
mode or active-constraint switches a two-sided derivative can fail.

The information comparison uses this same physical model and objective:
one observer sees E and A+B at each section, another sees each lane.
Admit model-aware inference from allowed histories; choose advice before
the subsequent revealing discharge. A compact observer adds the final-A
density, sufficient on the existing two-state diagnostic. Separately,
compare gradients with bounded scalar search using the same information,
and retain maximum advice as a reference. These definitions isolate what
information is needed from how the advice is optimized.
The \href{vsl_joined_candidate_contract_and_qualification.pdf}{companion formulation}
specifies all bias branches and finite-domain conventions.

\par\pdfsavepos
\write16{REVIEW-MAIN-END-PAGE: \thepage; Y-SP: \the\pdflastypos}
\makeatletter
\renewcommand{\section}{\@ifstar{\head}{\head}}
\makeatother
\begin{thebibliography}{8}
\setlength{\itemsep}{4pt}\setlength{\parskip}{0pt}
\bibitem{jinjin} H.-Y. Jin and W.-L. Jin, ``Control of a lane-drop bottleneck
through variable speed limits,'' arXiv:1310.2658, 2013 (archived manuscript
dated Oct. 29, 2018). [Online]. Available:
\url{https://arxiv.org/abs/1310.2658}.
\bibitem{wei} L. Wei, Y. Han, and M. Wang, ``A first-order link-based flow
model with variable speed limits and capacity drops for freeway networks,''
\emph{Comput.-Aided Civ. Infrastruct. Eng.}, vol.~40, no.~21,
pp.~3200--3217, 2025, doi:
\href{https://doi.org/10.1111/mice.13492}{10.1111/mice.13492}.
\bibitem{ns} H.~H.~S. Nagalur Subraveti, V.~L. Knoop, and B. van Arem,
``First order multi-lane traffic flow model---an incentive based macroscopic
model to represent lane change dynamics,'' \emph{Transportmetrica B, Transp.
Dyn.}, vol.~7, no.~1, pp.~1758--1779, 2019 (archived article header: 2020),
doi: \href{https://doi.org/10.1080/21680566.2019.1700846}{10.1080/21680566.2019.1700846}.
\bibitem{han16} Y. Han, Y. Yuan, A. Hegyi, and S.~P. Hoogendoorn,
``New extended discrete first-order model to reproduce propagation of jam
waves,'' \emph{Transp. Res. Rec.}, vol.~2560, pp.~108--118, 2016, doi:
\href{https://doi.org/10.3141/2560-12}{10.3141/2560-12}.
\bibitem{han17} Y. Han, A. Hegyi, Y. Yuan, S. Hoogendoorn,
M. Papageorgiou, and C. Roncoli, ``Resolving freeway jam waves by discrete
first-order model-based predictive control of variable speed limits,''
\emph{Transp. Res. C, Emerg. Technol.}, vol.~77, pp.~405--420, 2017, doi:
\href{https://doi.org/10.1016/j.trc.2017.02.009}{10.1016/j.trc.2017.02.009}.
\bibitem{carey} M. Carey, C. Balijepalli, and D. Watling, ``Extending the
cell transmission model to multiple lanes and lane-changing,''
\emph{Netw. Spat. Econ.}, vol.~15, pp.~507--535, 2015, doi:
\href{https://doi.org/10.1007/s11067-013-9193-7}{10.1007/s11067-013-9193-7}.
\bibitem{thesis} H.~H.~S. Nagalur Subraveti, \emph{Lane-specific Traffic
Flow Control}, Ph.D. dissertation, Delft Univ. Technol., Delft, The
Netherlands, 2021, ch.~4, doi:
\href{https://doi.org/10.4233/uuid:e405a549-b1fe-4fa1-8090-4429612359d9}{10.4233/uuid:e405a549-b1fe-4fa1-8090-4429612359d9}.
\bibitem{wright} M.~A. Wright, G. Gomes, R. Horowitz, and
A.~A. Kurzhanskiy, ``On node models for high-dimensional road networks,''
\emph{Transp. Res. B, Methodol.}, vol.~105, pp.~212--234, 2017, doi:
\href{https://doi.org/10.1016/j.trb.2017.09.001}{10.1016/j.trb.2017.09.001}.
\end{thebibliography}
\end{document}
